\section{Structure of the proof}\label{sec:structure}

We work throughout this section in $\ZFm$.

To prove Theorem~\ref{thm:main}, it suffices to show that every set can
be well-ordered. We shall place an arbitrary set inside an algebraic
field extension and obtain its well-ordering from the following
theorem. This is the first reduction of the proof: it turns the
set-theoretic question into a question about extending a well-ordering
from a field to a larger field.

\begin{overviewtheorem}\label{thm:algebraic}
Assume $\MC$. Let $L/K$ be a normal separable algebraic extension.
Assume that $\prec$ is a well-ordering of $K$ and that, for every well-orderable
field $F$ with $K\subseteq F\subseteq L$, the $F$-vector space $L$ has a basis.
Then $L$ is well-orderable by a well-ordering in which $(K, \prec)$ is an initial segment.
\end{overviewtheorem}

To apply this theorem, we need a suitable extension containing a copy
of the set we started with. The next theorem provides it using
$\MC$, which is already available from Blass's theorem.

\begin{overviewtheorem}\label{prop:embedding}
Assume $\MC$. For every set $X$ there exist a well-orderable field $K$
of characteristic zero, a normal separable algebraic extension $L/K$,
and an injection $X\to L$.
\end{overviewtheorem}

The construction uses L\'evy's partition lemma to write $X$ as an
ordinal-indexed union of disjoint finite sets. Replacing the elements
of $X$ by indeterminates gives a rational function field $L$. For each
finite set in the partition, we form the polynomial whose roots are its
indeterminates, and let $K$ be generated by the coefficients of these
polynomials. The coefficients have a well-ordered indexing, which
makes $K$ well-orderable, while the polynomials make $L/K$ algebraic,
normal and separable. The details are given in
Section~\ref{sec:main-proof}.

\begin{proof}[Proof of Theorem~\ref{thm:main}]
By Theorem~\ref{thm:blass}, $\MC$ holds. Let $X$ be an arbitrary set.
By Theorem~\ref{prop:embedding}, there exist a well-orderable
field $K$ of characteristic zero, a normal separable algebraic
extension $L/K$, and an injection $X\to L$.
Fix a well-ordering of $K$. Every intermediate field $F$ containing
$K$ has characteristic zero, so the hypothesis of
Theorem~\ref{thm:main} gives a basis of the $F$-vector space $L$.
Theorem~\ref{thm:algebraic} therefore gives a well-ordering of $L$,
which induces a well-ordering of $X$. Since $X$ was arbitrary,
$\AC$ follows.
\end{proof}

We now work backward from the conclusion of
Theorem~\ref{thm:algebraic}. If we can reach a well-ordered
intermediate field $F$ such that $L/F$ is finite, then a finite basis
gives a well-ordering of $L$ through its coordinates over $F$.
It is therefore enough to have a procedure that extends the current
well-ordered field, with the guarantee that whenever the procedure
adds nothing, the remaining extension is finite. The following theorem
provides precisely this step.

\begin{samepage}
\begin{overviewtheorem}\label{lem:adjoin}
Let $L/F$ be a normal separable algebraic extension, let $\prec$ be a
well-ordering of $F$, and let $\Ccal$ be a finite nonempty collection
of $F$-bases of $L$. From these data one can define an intermediate
field $F^+$, with $F\subseteq F^+\subseteq L$, and a well-ordering
$\prec^+$ of $F^+$ such that $(F,\prec)$ is an initial segment of
$(F^+,\prec^+)$ and
\[
 F^+=F\quad\Longrightarrow\quad [L:F]<\omega.
\]
\end{overviewtheorem}
\end{samepage}

Here the extension of the well-ordering is determined by the given
data. To prove Theorem~\ref{thm:algebraic} using this construction, we first apply
$\MC$ to the sets of bases over all well-orderable intermediate
fields, obtaining a finite nonempty family $\Ccal_F$ for each such
$F$. With these families fixed, we start at $K$ and repeatedly apply
the theorem. At limit stages we take unions. The requirement that each
order extend the previous one as an initial segment ensures that
these unions are still well-ordered.

The process must reach a stage at which the field no longer grows.
Indeed, let $\kappa=h(\Pow(L))$ be the Hartogs number of $\Pow(L)$,
the least ordinal that does not inject into the set of subsets of
$L$ (see \cite[Lemma~3.4, p.~29]{Jech2003}). A strictly increasing chain of intermediate fields of length
$\kappa$ would contradict its definition. At a stationary stage,
Theorem~\ref{lem:adjoin} gives a finite remaining extension, so the
finite-basis argument completes the well-ordering of $L$. This is
the proof of Theorem~\ref{thm:algebraic} from
Theorem~\ref{lem:adjoin}. Its formal details appear in
Section~\ref{sec:extension}.

It remains to explain the enlargement procedure in
Theorem~\ref{lem:adjoin}. Its definition is guided by a criterion for
the finiteness of $L/F$. To state that criterion, let
$\Ecal(L,F)$ denote the family of intermediate fields $E$ for which
$E/F$ is finite, normal and separable. If $B$ is an $F$-basis of $L$
and $E\in\Ecal(L,F)$, then $B\cap E$ is finite, since it is linearly
independent in the finite-dimensional $F$-vector space $E$. Put
\[
 q_{B,E}(x)=\prod_{b\in B\cap E}(x-b)\in E[x],
 \qquad
 P_{\Ccal,E}(z,x)=\prod_{B\in\Ccal}
                 \bigl(z-q_{B,E}(x)\bigr)\in E[z,x],
\]
where $\Ccal$ is a finite nonempty collection of $F$-bases of $L$ and
an empty product is $1$. The first polynomial records the finite
intersection with one basis. The second records the collection of
these polynomials without selecting a basis from $\Ccal$. Here $z$
is a second indeterminate, independent of $x$.

\begin{overviewtheorem}\label{lem:finite-bases}
Let $L/F$ be a normal separable algebraic extension and let $\Ccal$
be a finite nonempty family of $F$-bases of $L$. Suppose that, for every $E\in\Ecal(L,F)$,
\begin{equation*}
 P_{\Ccal,E}(z,x)\in F[z,x].
\end{equation*}
Then $L/F$ is finite.
\end{overviewtheorem}

This criterion suggests which elements should be adjoined in
Theorem~\ref{lem:adjoin}. Writing $\coeff(P)$ for the finite set of
coefficients of a polynomial $P$, we put
\[
 F^+=F\left(\bigcup_{E\in\Ecal(L,F)}\coeff(P_{\Ccal,E})\right).
\]
If $F^+=F$, all the polynomials $P_{\Ccal,E}$ have their coefficients
in $F$, and Theorem~\ref{lem:finite-bases} gives $[L:F]<\omega$.
To complete the enlargement step, we must also extend the given
well-ordering to $F^+$. Each $E\in\Ecal(L,F)$ is generated by the
roots of a polynomial over $F$. We can therefore index the
coefficients used above by polynomials in $F[x]$ and by pairs of
natural numbers recording the powers of $z$ and $x$. Since $F$ is
well-ordered, this is a well-ordered index set. Ordering the formal
expressions in these generators then gives the required extension
of the order, as shown in Section~\ref{sec:extension}.

The proof of Theorem~\ref{lem:finite-bases} is where Galois theory
enters. If $P_{\Ccal,E}$ has coefficients in $F$, every automorphism
of $E$ over $F$ sends a set $B\cap E$, with $B\in\Ccal$, to another
set $B'\cap E$ with $B'\in\Ccal$. This bounds the number of possible
conjugates of the coefficients of the polynomials $q_{B,E}$.
For one fixed basis in $\Ccal$, the resulting degree bound, together
with linear independence, forces that basis to be finite.
Section~\ref{sec:finite} proves this criterion using the elementary
Galois-theoretic facts established in Section~\ref{sec:galois}.
